\documentclass[11pt,a4paper]{article}

\usepackage[margin=2.4cm]{geometry}
\usepackage{amsmath,amssymb}
\usepackage{booktabs}
\usepackage{longtable}
\usepackage{listings}
\usepackage{xcolor}
\usepackage[colorlinks=true,linkcolor=blue!50!black,urlcolor=blue!50!black]{hyperref}
\usepackage{microtype}

\lstset{
  basicstyle=\ttfamily\small,
  breaklines=true,
  frame=single,
  framerule=0.3pt,
  rulecolor=\color{gray!50},
  backgroundcolor=\color{gray!5},
  columns=fullflexible,
  showstringspaces=false,
}

\newcommand{\code}[1]{\texttt{#1}}
\newcommand{\file}[1]{\texttt{\textbf{#1}}}

\title{\textbf{Worldline QMC: Code Structure}\\[0.3em]
       \large What each Python module does and why}
\author{Implementation notes}
\date{\today}

\begin{document}
\maketitle

\begin{abstract}
A file-by-file account of the Python in \code{quantum\_monte\_carlo/}. The
package \code{qmc/} is a conversion layer wrapped around the DSQSS/DLA worldline
Monte Carlo engine (vendored in \code{external\_dsqss/}, GPLv3, with one local
patch); the top-level scripts are the command-line entry point, the correctness
harness, and the benchmarks against the Chebyshev typicality code. The physics
and the estimator corrections are documented separately in
\code{qmc\_implementation.tex}; this document is about the software.
\end{abstract}

\tableofcontents

\newpage

%======================================================================
\section{Overview}

\subsection{Layout}

\begin{lstlisting}
quantum_monte_carlo/
|-- qmc/                       the package: everything reusable
|   |-- __init__.py            re-exports the submodules
|   |-- lattice.py             lattice definition; single source of truth
|   |-- parse.py               readers/writers for DSQSS's own file formats
|   |-- correlations.py        estimator algebra: DSQSS output -> physics
|   |-- storage.py             HDF5 writer in the Chebyshev schema
|   |-- console.py             terminal styling for the run log
|   `-- run.py                 orchestration; ties all of the above together
|-- run_qmc.py                 command-line entry point
|-- validate_ed.py             correctness harness (exact diagonalization)
|-- benchmark_vs_chebyshev.py  benchmark at zero field, three sites
|-- benchmark_field.py         benchmark at finite field, all four components
`-- plot_benchmark.py          figures for the zero-field benchmark
\end{lstlisting}

\subsection{Dependency structure}

The package is a shallow stack with no cycles. Only \code{run.py} knows about
more than one concern:

\begin{center}
\begin{tabular}{ll}
\toprule
Module & Imports from the package \\
\midrule
\code{lattice.py}     & --- (numpy; \code{h5py} lazily, inside one method) \\
\code{parse.py}       & --- \\
\code{correlations.py}& --- \\
\code{storage.py}     & --- \\
\code{console.py}     & --- \\
\code{run.py}         & all five of the above \\
\bottomrule
\end{tabular}
\end{center}

The four leaf modules are independently testable and contain no I/O beyond their
own narrow file formats. That matters because the estimator corrections
(Section~\ref{sec:correlations}) are the part most likely to need revisiting,
and they live in a module with no dependencies at all.

\subsection{Where the work actually happens}

A single run passes through the package in this order:

\begin{enumerate}
  \item \code{lattice.build} turns a spec string such as \code{square:6x4} into a
    \code{PeriodicLattice}.
  \item \code{run.run} resolves parameters and calls \code{run.run\_dsqss}.
  \item \code{lattice.std\_toml} renders the DSQSS input; \code{dla\_pre}
    expands it into DSQSS's XML files.
  \item \code{parse.write\_full\_bz\_wavevectors} overwrites one of those files
    (\code{wv.xml}) to fix a Brillouin-zone coverage bug.
  \item \code{dla} runs; its progress is streamed through
    \code{run.\_echo\_progress} and \code{console.Style}.
  \item \code{parse.read\_results} and \code{parse.read\_displacement\_kinds}
    read the raw output.
  \item \code{correlations.*} converts the raw estimators into
    $C^{xx}$, $C^{xy}$, $C^{zz}$.
  \item \code{storage.store} writes the HDF5 file in the Chebyshev schema.
\end{enumerate}

%======================================================================
\section{The package}

%----------------------------------------------------------------------
\subsection{\file{qmc/\_\_init\_\_.py} --- the package surface}

Five lines. It re-exports the six submodules so that \code{import qmc} followed
by \code{qmc.lattice.build(...)} works, and declares \code{\_\_all\_\_}. It
contains no logic deliberately: importing the package should not have side
effects, and in particular should not require the DSQSS binaries to be present.
Scripts import what they need directly
(\code{from qmc import run as run\_mod}), so the re-export is a convenience
rather than the expected entry route.

%----------------------------------------------------------------------
\subsection{\file{qmc/lattice.py} --- geometry and the single source of truth}

\paragraph{Purpose.} Defines the lattice and, critically, defines it \emph{once}
for both codes.

\paragraph{Why it is written this way.} The Chebyshev code stores lattices as
raw $J_{ij}$ matrices with an arbitrary site ordering; some were built by hand
with explicitly listed periodic pairs, and carry no geometric information at
all. Recovering the Bravais structure from such a matrix in order to line up
site indices with a QMC run is a graph-isomorphism problem in general. This
module avoids that entirely by inverting the direction of information flow: one
\code{PeriodicLattice} object emits \emph{both} the DSQSS input and the
Chebyshev-format coupling file, so the two codes agree on what ``site 5'' means
by construction rather than by inspection.

\paragraph{Class \code{PeriodicLattice}.}

\begin{longtable}{@{}p{0.30\textwidth}p{0.64\textwidth}@{}}
\toprule
Member & Role \\
\midrule
\endhead
\code{\_\_init\_\_(size, J, name)} &
  Stores \code{size} (per-dimension lengths), coupling \code{J}, and precomputes
  \code{coords}, the coordinate of every flat site index. \\
\code{\_default\_name()} &
  Produces \code{Chain\_NN\_PBC\_N=8}, \code{Square\_NN\_PBC\_N=24} etc.,
  matching the Chebyshev naming so output filenames line up. \\
\code{\_index\_to\_coord(i)} &
  Row-major, x-fastest unpacking, deliberately identical to DSQSS's
  \code{index2coord}: $i = x + L_x y + L_x L_y z$. \\
\code{\_coord\_to\_index(c)} & The inverse. \\
\code{coupling\_matrix()} &
  The $J_{ij}$ matrix. Contributions are \emph{accumulated} rather than
  assigned, which reproduces DSQSS's bond multiplicity: along a periodic
  direction of length 2 the $+1$ and $-1$ neighbours coincide and DSQSS lays
  down two bonds, giving an effective $2J$. \\
\code{displacement\_vector(i,j)} &
  Minimum-image displacement, using the same wrapping convention as
  \code{dsqss.displacement}. Used to identify correlation bins. \\
\code{is\_bipartite} (property) &
  True when every side length above 2 is even. Gates the antiferromagnet, which
  otherwise has a sign problem. \\
\code{marshall\_sign(i,j)} &
  $(-1)^{\sum_d r_d}$. Undoes the sublattice gauge DSQSS uses to make the
  antiferromagnet sign-free; applied to transverse components only, and only
  for $J>0$. \\
\code{write\_coupling\_file(path)} &
  Writes \code{all/J\_ij} plus the \code{num\_Spins} attribute --- the
  Chebyshev \code{Couplings/} format. \code{h5py} is imported inside the method
  so the module stays importable without it. \\
\code{std\_toml(...)} &
  Renders the DSQSS \code{std.toml}. Applies the Hamiltonian mapping
  $J_z = J_{xy} = -J$, $h = -h_z$. \\
\bottomrule
\end{longtable}

\paragraph{Function \code{build(spec, J)}.} Parses \code{chain:L},
\code{square:LxL}, \code{cube:LxLxL}. A single number expands over all
dimensions (\code{square:4} $\to$ \code{4x4}); explicit forms may be anisotropic
(\code{square:8x4}). Raises a clear \code{ValueError} on an unknown kind or a
dimension-count mismatch, which \code{run\_qmc.py} turns into a one-line CLI
error.

%----------------------------------------------------------------------
\subsection{\file{qmc/parse.py} --- DSQSS's own file formats}

\paragraph{Purpose.} Everything that reads or writes a DSQSS-native file. This
is the module that absorbs the engine's idiosyncrasies so nothing else has to.

\begin{longtable}{@{}p{0.34\textwidth}p{0.60\textwidth}@{}}
\toprule
Function & Role \\
\midrule
\endhead
\code{read\_results(path)} &
  Parses \code{R <name> = <mean> <error>} lines, the format DSQSS uses for every
  observable in every output file. Returns \code{\{name: (mean, error)\}}. \\
\code{read\_displacement\_kinds(path)} &
  Reads \code{disp.xml}, mapping each ordered site pair to its displacement bin.
  Uses a regular expression rather than an XML parser because the file is not
  well-formed XML --- it has bare comments and repeated top-level tags. \\
\code{wavevector\_list(size)} &
  The \emph{full} Brillouin zone as integer $k$ indices. \\
\code{write\_full\_bz\_wavevectors(...)} &
  Writes a \code{wv.xml} covering every $k$, replacing the one \code{dla\_pre}
  generates. \\
\code{collect\_tau\_series(...)} &
  Gathers \code{<prefix><kind>t<itau>} entries into \code{(nkinds, ntau)} arrays.
  Missing entries become NaN rather than zero, so a truncated or partial output
  file is visible instead of silently producing plausible-looking numbers. \\
\bottomrule
\end{longtable}

\paragraph{Why the full Brillouin zone.} DSQSS's own generator emits the product
set $k_d \in \{0,\dots,L_d/2\}$ and relies on $S(k) = S(-k)$ to cover the rest.
That is a valid fundamental domain for $k \to -k$ in one dimension but
\emph{not} in two or more: on a $4\times4$ lattice the orbit
$\{(1,3),(3,1)\}$ lies entirely outside it, so those $k$ are never measured and
the back-transform to real space silently loses their contribution. Measuring
the whole zone makes the transform exact with uniform weight and removes the
multiplicity bookkeeping. The bug this fixes produced a $261\sigma$ violation of
an exact symmetry; see \code{qmc\_implementation.tex}.

%----------------------------------------------------------------------
\subsection{\file{qmc/correlations.py} --- estimator algebra}
\label{sec:correlations}

\paragraph{Purpose.} Converts what DSQSS measures into what is physically
wanted. This is the smallest module and the most subtle; every function here
exists because the naive reading of the DSQSS output is wrong in some specific
way.

\begin{longtable}{@{}p{0.34\textwidth}p{0.60\textwidth}@{}}
\toprule
Object & Role \\
\midrule
\endhead
\code{CLASS\_C\_ROWS} &
  \code{("xx","xy","yx","zz")}. The row order every output file uses, matching
  symmetry class C of the Chebyshev \code{CorrelationTensor}. \\
\code{transverse\_from\_cf(...)} &
  $C^{xx} = \code{cf}/2$. The stock worm estimator does not record whether the
  head is $S^+$ or $S^-$, so what it reports is the symmetrised combination
  $\tfrac12[G^{+-}+G^{-+}]$, which is exactly $2C^{xx}$. \\
\code{\_reverse\_tau(a)} &
  Maps $i \mapsto (n_\tau - i) \bmod n_\tau$. \\
\code{xy\_from\_antisymmetric(...)} &
  $\mathrm{Im}\,C^{xy} = \tfrac12 \mathcal{R}\,\code{af}$, where \code{af} is the
  signed histogram added by the engine patch. The $\tau$ reversal is required
  because the worm bins the head--tail time displacement with the opposite
  orientation; that is invisible in the symmetric channel (even about
  $\beta/2$) but negates this one. \\
\code{zz\_from\_structure\_factor(...)} &
  Inverse Fourier transform of $S^{zz}(k,\tau)$ to real space. Validates that
  the $k$ grid spans the full zone, checks the input for finiteness, performs
  the transform under a suppressed-warning block, then checks the output again. \\
\code{mirror\_to\_full\_beta(...)} &
  Extends a $[0,\beta)$ series to $\tau=\beta$, with a sign flip for
  antisymmetric components. \\
\code{build\_tensor(...)} &
  Allocates the zero-filled \code{(nsites, nrows, ntau)} arrays. \\
\bottomrule
\end{longtable}

\paragraph{The suppressed warnings.} numpy built against Apple's Accelerate BLAS
leaves floating-point exception flags set inside its vectorised kernels, and
numpy reports them after the call even when every input and output is finite ---
confirmed by feeding finite random input, where the warning fires while the
result still matches \code{einsum} to $7\times10^{-15}$. Because the inputs are
checked for finiteness immediately before and the outputs immediately after, the
\code{errstate} block hides a known false alarm rather than a real condition.

%----------------------------------------------------------------------
\subsection{\file{qmc/storage.py} --- the output file}

\paragraph{Purpose.} Writes HDF5 in exactly the Chebyshev schema, so existing
\code{Plot\_*.py} scripts read QMC output unchanged.

\begin{longtable}{@{}p{0.34\textwidth}p{0.60\textwidth}@{}}
\toprule
Function & Role \\
\midrule
\endhead
\code{build\_filename(...)} &
  Reproduces \code{Storage\_Concept::create\_file}'s naming rule. Optional
  fragments appear only when they differ from their defaults. \code{J} is an
  addition to the C++ scheme: there the couplings live in a named source file,
  whereas here $J$ is a run parameter, and without it a ferromagnetic run would
  silently overwrite an antiferromagnetic one at the same $\beta$. \\
\code{\_fmt(x)} &
  Formats like a C++ \code{ostream}: integers without a trailing \code{.0}, so
  \code{beta=2} not \code{beta=2.0}. \\
\code{MAX\_TRIES}, \code{resolve\_path(...)} &
  Collision avoidance, following the C++ writer: append \code{X} to the stem
  until a free name is found, up to four. Once all five variants exist the last
  is overwritten, so a repeated scan cannot fill the disk. (The C++ writer
  raises at that point instead.) \\
\code{store(...)} &
  Writes \code{/parameters} as attributes, the four
  \code{/results/*/<site>-0} datasets each with the same \code{info} attribute
  string as the C++ writer, and \code{/runtime\_data}. Validates that the number
  of sites matches the array shape. \\
\bottomrule
\end{longtable}

Every run writes four rows regardless of field. Both the transverse and
longitudinal channels are measured on every run, so collapsing to a single row
at $h_z=0$ would discard the independent $C^{zz}$ estimate --- the cheapest
available check on $C^{xx}$, since isotropy requires the two to agree.

%----------------------------------------------------------------------
\subsection{\file{qmc/console.py} --- terminal styling}

\paragraph{Purpose.} ANSI styling for the run log, and nothing else. Isolated in
its own module so that presentation never entangles with computation.

\code{supports\_colour(stream)} returns true only for a real terminal, and
respects \code{NO\_COLOR} and \code{TERM=dumb}. \code{Style} exposes semantic
methods --- \code{head}, \code{value}, \code{ok}, \code{warn}, \code{bad},
\code{path}, \code{bar} --- which become no-ops when colour is unavailable.

Two design points are deliberate. Style is applied \emph{after} padding, so
column alignment is byte-identical with and without colour; and there is no
faint/dim style, because it renders close to unreadable on terminals whose
palette already has low contrast. Ordinary text keeps the terminal's default
colour and emphasis is carried by bold and hue alone.

%----------------------------------------------------------------------
\subsection{\file{qmc/run.py} --- orchestration}

\paragraph{Purpose.} The only module that knows about all the others. It drives
the external binaries, converts their output, and reports progress.

\subsubsection*{Locating the engine}

\code{INSTALL\_DIR} defaults to \code{dsqss\_install/} beside the package and can
be overridden with the \code{DSQSS\_INSTALL} environment variable.
\code{\_binary(name)} raises a message naming the expected path rather than
letting a bare \code{FileNotFoundError} escape.

\subsubsection*{Running the sampler}

\code{run\_dsqss} writes \code{std.toml}, invokes \code{dla\_pre} to expand it,
overwrites \code{wv.xml} with the full Brillouin zone, then runs \code{dla}
(under \code{mpirun} when \code{ncores > 1}).

The subprocess is \emph{streamed} rather than captured, so progress appears
while the sampler runs instead of arriving in one lump at the end. The full log
is still retained in memory so that a non-zero exit can be reported with
context. \code{\_find\_sf\_output} locates the structure-factor file by reading
the name back out of \code{param.in}, rather than assuming it.

\code{\_SET\_DONE} and \code{\_NCYC} are the regular expressions that recognise
DSQSS's own progress chatter; \code{\_echo\_progress} reformats a matching line
into a progress bar and returns whether it recognised it.
\code{\_format\_seconds} renders durations as \code{45s}, \code{2m31s},
\code{1h04m}.

\subsubsection*{Converting}

\code{convert} is the heart of the module. It reads the displacement bins and
both raw estimator files, applies the three corrections, and assembles the
\code{(nsites, 4, ntau)} tensors. It detects whether the engine is patched by
checking for the \code{A} entries; an unpatched build leaves $C^{xy}$ at zero
rather than producing something wrong. A structure factor that is partially
populated raises immediately, since that would mean the full Brillouin zone was
not measured.

\subsubsection*{Reporting}

\code{print\_banner} lists the parameters before sampling starts. It
deliberately does \emph{not} name the destination file, because the name is only
settled at write time. \code{\_warn\_nonfinite\_errors} flags error bars that
came back NaN --- DSQSS forms the error as
$\sqrt{\langle x^2\rangle - \langle x\rangle^2}$, which with few sets on a barely
fluctuating observable can go slightly negative through cancellation. The
correlations are unaffected, so this warns rather than aborting.
\code{\_print\_destination} reports the final path and any \code{X} suffix.

\subsubsection*{The entry point}

\code{run(...)} is the function every script calls. It builds the lattice,
validates the site list, rejects a frustrated antiferromagnet up front, runs,
converts, resolves the output path, writes, and cleans up the scratch directory
unless \code{keep\_workdir} is set.

%======================================================================
\section{Top-level scripts}

\subsection{\file{run\_qmc.py} --- command-line entry point}

A thin argparse wrapper over \code{run.run}, with options deliberately mirroring
the Chebyshev executable (\code{--beta}, \code{--num\_TimePoints}, sites,
cores). It adds \code{--write-couplings} to emit the matching coupling file, and
\code{--quiet} to reduce output to the bare path so that
\code{path=\$(run\_qmc.py ... --quiet)} works in scripts.

Configuration errors (\code{ValueError}) are caught and printed as a single line
with exit code 2; a traceback would add nothing for a mistyped lattice spec or
an out-of-range site.

\subsection{\file{validate\_ed.py} --- the correctness harness}

Builds the Hamiltonian
$H = \sum_{i<j} J_{ij}\,\mathbf{S}_i\!\cdot\!\mathbf{S}_j + h_z\sum_i S^z_i$
independently, diagonalises it densely, and evaluates
\begin{equation*}
  C^{ab}(\tau) = \frac{1}{Z}\sum_{n,m} e^{-\beta E_n} e^{\tau(E_n - E_m)}
                 A_{nm} B_{mn},
\end{equation*}
then compares every component of the QMC output against it. This is the
ground-truth test: it shares no code with the QMC path beyond the lattice
definition.

Agreement is accepted if the deviation is within tolerance in units of $\sigma$
\emph{or} below a small absolute threshold, because $\tau=0$ has a vanishing
statistical error --- it is fixed by an operator identity --- and a pure
$\sigma$ test is meaningless there.

Two constraints are enforced up front: at most 14 sites (dense diagonalisation),
and a bipartite lattice unless $J<0$. The help text steers towards
\code{square:3x3 --J=-1} when the $C^{zz}$ path is being touched, because
lattices with short sides make $k \to -k$ trivial and hide Brillouin-zone
coverage bugs --- an earlier suite of \code{square:4x2} and \code{cube:2x2x2}
passed while $C^{zz}$ was wrong.

\subsection{\file{benchmark\_vs\_chebyshev.py} --- zero-field benchmark}

Compares against real Chebyshev output on the 24-site periodic square, at a size
where exact diagonalisation is impossible ($2^{24}$ states). The lattice is
\code{square:6x4}, whose coupling matrix is byte-identical to
\code{Couplings/Square\_NN\_PBC\_N=24.hdf5}, so sites 0, 1 and 7 are the local,
nearest-neighbour and next-nearest correlations.

The two codes use different $\tau$ conventions: Chebyshev stores
\code{linspace(0, beta/2, num\_TimePoints)} --- its \code{Tmax} is $\beta/2$ ---
while QMC stores a uniform grid on $[0,\beta)$. The comparison is made at the
QMC grid points inside $[0,\beta/2]$ with the dense reference interpolated onto
them, so the noisy data is never interpolated. $\tau=0$ is excluded from the
pull statistic and reported separately as an absolute deviation.

\code{--skip-existing} reuses completed runs, so the comparison can be re-run
without repeating the sampling.

\subsection{\file{benchmark\_field.py} --- finite-field benchmark}

The same lattice at $h_z \neq 0$. The field data is local only (site 0) but
carries all four rows, so this exercises what vanishes at zero field:
$C^{xx}$ and $C^{zz}$ separately, and $\mathrm{Im}\,C^{xy}$ --- which is
measurable only with the patched estimator and cannot be checked by exact
diagonalisation at this size.

Three file-format quirks are handled. The field files predate the per-site group
layout, so \code{results/Re\_correlation} is a flat $(4, n_\tau)$ dataset rather
than a group. Error datasets are named \code{Re\_stds} in newer files and
\code{Re\_stddev} in older ones. The $h_z=2$ set carries \emph{no} uncertainties
at all --- those runs are reported as absolute agreement only, in a separate
block, rather than inventing a reference error. \code{AVAILABLE} records which
$(h_z, J)$ combinations exist and at which $\beta$.

\subsection{\file{plot\_benchmark.py} --- figures}

Produces one figure per coupling sign: correlations on top, pull underneath, one
column per site, all inverse temperatures overlaid. Plotted against $\tau/\beta$
so every temperature shares the horizontal range $[0,\tfrac12]$, which makes the
approach to the symmetry point at $\beta/2$ line up across temperatures and lets
the pull traces be read as flat bands. $\tau=0$ is dropped from the pull panel
for the same reason it is dropped from the statistic.

%======================================================================
\section{Conventions that recur}

Several habits appear throughout and are worth stating once.

\paragraph{Missing data becomes NaN, never zero.} \code{collect\_tau\_series}
fills absent entries with NaN so that a truncated output file is visible.
Several call sites then check explicitly and raise.

\paragraph{Checks bracket suppressed warnings.} Where a warning is silenced, the
inputs are validated immediately before and the outputs immediately after, so
the suppression cannot hide a real condition.

\paragraph{Physics identities are used as tests.} $C^{xx}(0)=\tfrac14$,
$C^{xx}=C^{zz}$ at $h_z=0$, $C^{yx}=-C^{xy}$, and $C^{xy}(\beta/2)=0$ all hold
exactly and are checked rather than assumed. The isotropy identity in particular
is size-independent, and is what exposed the Brillouin-zone bug at a system size
where no exact reference exists.

\paragraph{Configuration errors are not tracebacks.} Anything the user can
mistype raises \code{ValueError} with a message naming the fix, which the CLI
prints as one line with exit code 2.

\paragraph{Presentation is separable.} Every reporting function takes
\code{out} and \code{style}, and \code{progress} defaults to \code{False} in
\code{run\_dsqss}, so library use is silent and the harnesses stay readable.

\end{document}
